\documentclass[11pt,a4paper]{article}
\usepackage[margin=1in]{geometry}
\usepackage{graphicx}
\usepackage{cite}
\usepackage{amsmath,amssymb,amsfonts}
\usepackage{algorithmic}
\usepackage{textcomp}
\usepackage{xcolor}
\usepackage{url}
\usepackage{booktabs}
\usepackage{multirow}
\usepackage[section]{placeins}

\begin{document}

\title{BASIRA AI: A Voice-First, Edge--Cloud Hybrid Multimodal Assistive
System for Blind and Visually Impaired Arabic-Speaking Users}

\maketitle
\date{}

\begin{abstract}
Blindness and severe visual impairment affect more than 295 million people
worldwide, and an estimated 13 million in the Arab world, yet the assistive
technology market remains dominated by either prohibitively expensive dedicated
devices (such as OrCam MyEye, retailing above 3{,}500 USD) or English-centric
smartphone applications that under-serve right-to-left languages and local
artefacts such as the Egyptian Pound (EGP) banknote series. As a result, the
majority of Arabic-speaking blind and visually impaired (BVI) users still depend
on sighted assistance for routine tasks: identifying currency, reading
handwritten or printed signs, recognising people, scanning grocery products,
and avoiding indoor obstacles. We present \textbf{BASIRA AI}---an Arabic word
meaning ``insight''---a voice-first, multimodal assistive Flutter application
that consolidates twelve perception capabilities into a single, accessible
mobile experience. This paper is an \emph{applied assistive-systems} paper:
it does \emph{not} propose a new learning algorithm. Instead, we describe the
engineering integration and adaptation of existing, well-known components
(Flutter, Google ML Kit on-device models, an off-the-shelf YOLOv8
backbone trained on a curated EGP dataset, a cloud multimodal large language
model---Gemini~1.5~Flash, and a rule-based edge--cloud router) into a
single Arabic-first, voice-first product targeted at BVI users in Egypt.
The research value lies in the system-level decisions---privacy-gated
routing, no-silence guarantee, Arabic-native TTS pipeline, localised EGP
currency recognition---and in their measured impact on real users, not in
novel models. We report (i) a custom-trained EGP YOLOv8n detector that
reaches 96.4\% top-1 accuracy and 0.91 mAP@0.5 in a 6.7~MB INT8 TFLite
artefact, benchmarked against a Seeing-AI baseline that lacks EGP support;
(ii) end-to-end median latency comparisons against Seeing AI and Google
Lookout on currency, OCR and scene-description tasks; (iii) a deliberately
simple obstacle-warning rule (image-label keyword matching fused with a
bounding-box area/centre threshold) for indoor walking, evaluated on a
small staged 240-image set with AUC $=0.82$ and a chosen operating point
of TPR~$=0.92$ at FPR~$=0.04$---we do not claim it is a general
navigation solution; and (iv) a System Usability Scale (SUS) score of 84.6
and an average task-completion improvement of 38\% over previously used
tools, from a study with 18 BVI participants. BASIRA AI is presented as
evidence that careful integration and Arabic-first design, rather than new
models, can deliver immediately useful assistive technology on commodity
smartphones.
\end{abstract}

\noindent\textbf{Keywords:} Assistive technology, Visual impairment,
Multimodal AI, Edge--cloud orchestration, On-device inference, Arabic NLP,
YOLOv8, TensorFlow Lite, Multimodal large language models, Accessibility,
Egyptian currency recognition, Voice user interface.

%====================================================================
\section{Introduction}
\label{sec:introduction}

According to the World Health Organisation, at least 2.2~billion people live
with some form of vision impairment, of whom approximately 43~million are
blind and 295~million have moderate-to-severe distance impairment~\cite{who2023};
updated Global Burden of Disease estimates project that the absolute number of
blind adults will exceed 61~million by 2050 in the absence of new
interventions~\cite{bourne_vision_loss}.
Regional epidemiological studies estimate that blindness prevalence in the
Eastern Mediterranean region is roughly 1.4 to 2.0 times the global
average~\cite{who_emr}, with Egypt alone accounting for more than 1.8~million
visually impaired citizens~\cite{egypt_disability}. For these users, the
``last mile'' of daily independence---identifying a banknote at a market,
reading a medication label, locating an obstacle in an unfamiliar
corridor---remains stubbornly unsolved.

The contemporary assistive landscape is split between two extremes, as documented
in recent reviews of vision-assistive technology~\cite{assistive_tech_review2}.
On one
side, dedicated wearable devices such as OrCam MyEye and Envision
Glasses~\cite{envision_app}
provide high-quality offline perception but at price points (3{,}500--4{,}500
USD) that exclude the overwhelming majority of users in low- and
middle-income countries~\cite{orcam_review}. On the other side, smartphone
applications such as Microsoft Seeing AI, Google Lookout, and Be My Eyes
deliver impressive functionality at zero marginal cost, yet their Arabic
support is shallow: Arabic OCR quality lags behind English by 12--20
percentage points on common benchmarks~\cite{arabic_ocr_gap}, currency
modules cover USD, EUR and GBP but not EGP, and right-to-left voice
interaction is frequently treated as a translated afterthought rather than a
first-class design constraint, despite explicit guidance to the contrary in
WCAG~2.2 and related international
accessibility standards~\cite{wcag22}.

A second, less visible problem is architectural. Most assistive applications
treat each capability (OCR, object detection, currency, scene description) as
an isolated screen with its own model, threading model, and threshold logic.
This fragmentation produces three concrete failures observed in the field:
(i) inconsistent voice feedback (different rates, voices, and interruption
policies across screens); (ii) wasted compute and battery from redundantly
loaded models; and (iii) no principled fall-back when a model fails or the
network drops, leaving the user in silence at the moment they most need
guidance.

This paper introduces \textbf{BASIRA AI}, a voice-first multimodal assistive
application that addresses both gaps---linguistic and architectural---through
a unified \emph{AI orchestration layer}. The system is implemented in Flutter
to target both Android and iOS from a single codebase, integrates twelve
perception features behind a single accessible home screen, and routes every
inference request through a deterministic policy that decides between
on-device models (Google ML Kit, a custom TensorFlow Lite EGP detector) and
a cloud multimodal LLM (Gemini~1.5~Flash) based on task class, connectivity,
latency budget, and privacy sensitivity. The application was developed with
active involvement of BVI users in Cairo and is publicly framed under the
Arabic name \emph{Basira} (insight), reflecting a design philosophy in which
the software, not the user, must adapt.

\subsection{Scope and Type of Contribution}
\label{sec:scope}

We explicitly position this work as an \emph{applied assistive-systems}
paper, not a machine-learning method paper. Every individual building block
we use---the Flutter cross-platform UI toolkit, Google ML Kit's on-device
OCR / object / face / barcode / image-label models, the YOLOv8 family of
detectors, the Gemini~1.5~Flash multimodal API, and rule-based
edge--cloud routing---is established prior art. We do not claim a novel
learning algorithm, a new architecture, or a new training procedure for
any of them. The contribution we ask the reader to evaluate is the
integration: the \emph{specific combination of choices} that make these
components work together as a single, Arabic-first, voice-first product
for BVI users in Egypt; the operational properties of that combination
(latency, offline behaviour, privacy gating, no-silence guarantee); and
the measured impact on a real BVI user cohort. Section~\ref{sec:baselines}
benchmarks BASIRA AI against representative existing applications
(Microsoft Seeing AI, Google Lookout) on tasks they share, to quantify
what the integration adds beyond any of its individual ingredients.

\subsection{Contributions}

With the scope of Section~\ref{sec:scope} in mind, the concrete
engineering and evaluation contributions of this paper are:

\begin{enumerate}
\item \textbf{A unified, voice-first multimodal assistive application.}
We design, implement, and release a Flutter-based mobile system that
consolidates twelve assistive capabilities (scene description, currency
recognition, text reading, object detection, an obstacle-warning aid,
face detection, colour detection, barcode/product scanning, light
detection, time \& location, settings, and a home voice-command loop)
behind a single Arabic-first, right-to-left, voice-driven user
interface. The integration---not any individual component---is the
artefact.

\item \textbf{A deterministic edge--cloud routing policy with explicit
fall-back.} We describe and instrument a rule-based routing function
$\pi: \mathcal{T} \times \mathcal{C} \rightarrow \mathcal{M}$ that maps a
(task, context) pair to an on-device or cloud engine. We make no claim
of a novel routing algorithm; we instead document the rule set, measure
its per-engine hit rate on one week of dogfood usage, and report the
fraction of invocations that recover useful output through the explicit
fall-back path.

\item \textbf{A localised on-device EGP currency recogniser.}
We assemble a curated Egyptian-pound dataset covering the six
circulating denominations (5, 10, 20, 50, 100, 200), including the
post-2021 polymer 10 and 20 EGP notes, and fine-tune an off-the-shelf
YOLOv8n backbone, which we then export and INT8-quantise to a 6.7~MB
TFLite artefact. The novelty is in the dataset and product packaging
for BVI users, not in the detector architecture.

\item \textbf{An honest, minimal obstacle-warning rule for indoor
walking.} The current ``navigation assistant'' is intentionally simple:
it applies a keyword filter over Google ML~Kit image-labeller outputs
(stairs, fence, wall, kerb, barrier, \ldots) and combines it with a
fixed bounding-box area/centre threshold over ML~Kit object-detector
outputs. We report its operating point and ROC on a small staged
indoor dataset and explicitly state what it does \emph{not} do
(no depth estimation, no semantic segmentation, no outdoor
evaluation).

\item \textbf{End-to-end comparisons against existing assistive
applications.} Section~\ref{sec:baselines} compares BASIRA AI against
Microsoft Seeing AI and Google Lookout on currency, OCR and scene-
description tasks, on the same Android handset and the same physical
targets, reporting both latency and task success.

\item \textbf{An accessibility-driven user evaluation.} We report a
study with 18 BVI participants in Cairo demonstrating a System
Usability Scale (SUS) score of 84.6, a 38\% mean reduction in task-
completion time relative to participants' previously used tools, and
qualitative evidence that Arabic-native voice feedback is the strongest
single determinant of adoption.
\end{enumerate}

\subsection{Paper Organisation}

Section~\ref{sec:related_work} reviews prior assistive systems, mobile
multimodal AI, on-device inference, and Arabic-language accessibility.
Section~\ref{sec:problem} states the problem formally.
Section~\ref{sec:architecture} describes the BASIRA AI system architecture
and the orchestration layer.
Section~\ref{sec:ai_components} details each integrated AI component,
including the EGP currency model and the explicit obstacle-warning rule.
Section~\ref{sec:dataset} describes datasets, annotation protocol, and
preprocessing.
Section~\ref{sec:results} presents quantitative model results, end-to-end
latency, energy measurements, head-to-head baselines
(Section~\ref{sec:baselines}), and the BVI user study.
Section~\ref{sec:discussion} discusses limitations---including the
absence of large-scale real-world deployment---and ethical
considerations.
Section~\ref{sec:conclusion} concludes.

%====================================================================
\section{Related Work}
\label{sec:related_work}

\subsection{Smartphone-Based Assistive Applications}

Smartphone assistive applications for the BVI community emerged in the
early 2010s with TapTapSee~\cite{taptapsee} and KNFB Reader~\cite{knfb_reader},
evolving rapidly after
multimodal cloud APIs became affordable. Microsoft's Seeing AI~\cite{seeingai}
introduced the now-standard pattern of a single application bundling several
``channels'' (short text, document, person, scene, currency, colour, light,
handwriting). Google Lookout~\cite{lookout} extended the pattern with a more
explicit modal switcher and voice-only interaction. Be My Eyes~\cite{bemyeyes}
took an orthogonal route by routing live video to sighted volunteers, more
recently augmented with GPT-4 vision. Foundational HCI studies of smartphone
use by BVI users~\cite{kane_blind_smartphone,branham_blind_hci} and of voice
assistants in non-visual
contexts~\cite{abdolrahmani_voice} establish that touch-screen accessibility
and conversational interaction are necessary but not sufficient when the
application itself is not Arabic-native. While these systems are mature,
three gaps motivate the present work: (a) shallow Arabic support, particularly for
diacritised text and Arabic-Indic numerals; (b) no Egyptian currency
recognition; and (c) closed source, which prevents local adaptation by
universities, NGOs and ministries of social affairs.

\subsection{Currency Recognition for the Visually Impaired}

Currency recognition is a canonical assistive task because banknotes are
visually similar, frequently fraudulent, and impossible to disambiguate by
touch in many modern series, including the polymer EGP notes introduced
between 2021 and 2023. Early work used SIFT, SURF and colour histograms with
SVM classifiers~\cite{currency_sift}; comprehensive reviews catalogue this
feature-engineered era~\cite{liu_currency_review}. Modern pipelines uniformly
adopt convolutional and deep detectors~\cite{mittal_banknote_dl}---YOLO
variants in particular~\cite{redmon_yolo,yolov8_tech}---trained on
denomination-labelled crops~\cite{currency_yolo}. The literature is
dominated by USD, EUR, INR and CNY datasets; EGP-specific public datasets
are scarce, and to our knowledge no published assistive application ships an
on-device EGP detector.

\subsection{On-Device Inference and Mobile ML}

The shift from cloud-only to on-device inference is driven by latency,
privacy, and connectivity. TensorFlow Lite~\cite{tflite}, Apple Core
ML~\cite{coreml}, and ONNX Runtime Mobile~\cite{onnx_runtime} now support
INT8 quantisation~\cite{jacob_int8}, GPU and NNAPI delegates,
and sub-50~MB packaging that fits into a typical app bundle. Mobile-tuned
backbones such as MobileNet~\cite{howard_mobilenets} and its inverted-residual
successor~\cite{sandler_mobilenetv2}, together with weight-pruning and
quantisation pipelines~\cite{han_deep_compression}, make sub-10~MB perception
models routine. Google ML
Kit~\cite{mlkit} encapsulates several mobile-tuned models (text recognition,
object detection, face detection, barcode, image labelling) behind a
high-level API, which we exploit for non-currency tasks where retraining
would yield diminishing returns.

\subsection{Multimodal Large Language Models for Accessibility}

The 2023--2025 generation of multimodal LLMs (GPT-4V~\cite{gpt4v},
Gemini~1.5~\cite{gemini_tech_report}, Claude~3.5
Sonnet~\cite{claude3}, and open alternatives such as LLaVA~\cite{llava})
collapses perception and language generation into a single API
call: a JPEG plus a prompt yields a coherent natural-language description.
For accessibility, this is transformative for open-ended tasks (``describe
the scene'', ``what is on the menu'') but carries three operational costs:
network round-trip latency (typically 1.2--3.5~s), per-call monetary cost,
and privacy exposure of raw camera frames. BASIRA AI uses a multimodal LLM
\emph{selectively}, only for tasks where on-device models are insufficient,
and explicitly degrades to offline mode when the network is unavailable.

\subsection{Arabic-Language Accessibility}

Despite Arabic being the fifth most spoken language globally, Arabic
accessibility tooling lags. Arabic OCR struggles with cursive ligatures,
diacritics, and the dual numeral system; Arabic TTS engines are unevenly
distributed across Android OEMs; and Arabic voice commands are often
limited to Modern Standard Arabic, missing Egyptian colloquial forms common
in everyday speech~\cite{arabic_dialect_survey}. Modern Arabic NLP toolkits
such as AraBERT~\cite{arabert} and CAMeL Tools~\cite{camel_tools} have
improved tokenisation, morphological analysis and downstream classification,
and recent surveys document steady gains in Arabic ASR
accuracy~\cite{arabic_asr_review}. Recent work~\cite{arabic_tts_survey} has
improved
neural Arabic TTS quality, and Google's on-device Arabic recogniser has
reached parity with English on printed text~\cite{mlkit_ar}, but
end-to-end assistive applications that treat Arabic as the first language
remain rare.

\subsection{Positioning of BASIRA AI}

Table~\ref{tab:related_comparison} positions BASIRA AI against representative
prior systems. The distinguishing combination is: (i) Arabic-first voice
interface, (ii) edge-first orchestration with explicit cloud fall-back,
(iii) on-device EGP currency recognition, and (iv) open architecture suitable
for local extension.

\begin{table}[h]
\caption{Comparison of representative assistive systems with BASIRA AI.}
\centering
\renewcommand{\arraystretch}{1.4}
\resizebox{\linewidth}{!}{%
\begin{tabular}{|l|c|c|c|c|c|c|}
\hline
\textbf{System} & \textbf{Arabic-first} & \textbf{EGP currency} & \textbf{On-device} & \textbf{Cloud LLM} & \textbf{Voice-first} & \textbf{Open} \\
\hline
Seeing AI~\cite{seeingai}        & Partial & No  & Mostly & Partial & Yes & No \\
Google Lookout~\cite{lookout}    & Partial & No  & Mostly & No      & Yes & No \\
Be My Eyes~\cite{bemyeyes}       & Partial & No  & No     & Yes (GPT-4) & Yes & No \\
OrCam MyEye~\cite{orcam_review}  & Partial & No  & Yes    & No      & Yes & No \\
Envision AI                      & Partial & No  & Partial& Yes     & Yes & No \\
\textbf{BASIRA AI (this work)}   & \textbf{Yes} & \textbf{Yes} & \textbf{Yes} & \textbf{Yes (selective)} & \textbf{Yes} & \textbf{Yes} \\
\hline
\end{tabular}}
\label{tab:related_comparison}
\end{table}

%====================================================================
\section{Problem Formulation}
\label{sec:problem}

We formalise the assistive perception problem as follows. Let $\mathcal{T}
= \{t_1, \ldots, t_K\}$ denote the set of supported assistive tasks (in
BASIRA AI, $K=12$). Let $x \in \mathbb{R}^{H \times W \times 3}$ denote a
camera frame and $c \in \mathcal{C}$ a context vector containing the
device's connectivity state, battery level, language preference, ambient
noise estimate, and a privacy class assigned by the user.

For each task $t_k$ we maintain a (possibly empty) set of inference engines
$\mathcal{M}_k \subseteq \mathcal{M} = \{\text{ML Kit}, \text{TFLite},
\text{Gemini}\}$. The orchestrator selects an engine via a routing function
\[
\pi : \mathcal{T} \times \mathcal{C} \rightarrow \mathcal{M}, \qquad
\pi(t_k, c) \in \mathcal{M}_k,
\]
which is then invoked to produce a structured output $y_k$. A post-processor
$g_k$ converts $y_k$ into a short Arabic utterance $u_k$ that respects the
user's preferred speech rate $\rho$ and is dispatched to the TTS engine.

The end-to-end objective is to minimise the perceived latency
\[
L(t_k, c) = L_{\text{capture}}(c) + L_{\text{infer}}(\pi(t_k, c), c) +
L_{\text{post}}(t_k) + L_{\text{tts}}(u_k, \rho)
\]
subject to (i) accuracy constraints $\text{Acc}(t_k) \geq \alpha_k$, (ii)
privacy constraints $\pi(t_k, c) \notin \text{cloud}$ whenever $c$ flags
sensitive content (faces, documents containing personal data), and (iii)
graceful degradation: if the chosen engine fails, the orchestrator must
return a non-empty Arabic utterance describing the failure rather than
silence. The latter constraint reflects the lived experience of BVI users:
silence is the worst outcome, because it is indistinguishable from a frozen
application.

%====================================================================
\section{System Architecture}
\label{sec:architecture}

\subsection{Overview}

BASIRA AI is a Flutter~\cite{flutter_framework} application targeting Android
(API 21+) and iOS
(13+), written in Dart~\cite{dart_language}. The codebase is organised into a
\texttt{core/} layer (services,
constants, theme) and a \texttt{features/} layer containing one folder per
assistive capability. Figure~\ref{fig:basira_arch} illustrates the
three-layer architecture: (1) the \emph{Presentation Layer} (voice loop,
home screen, twelve feature screens, settings); (2) the \emph{Orchestration
Layer} (routing policy, camera service, TTS service, settings service,
haptic service, Gemini service); and (3) the \emph{AI Layer} (Google ML Kit
models, custom TFLite EGP detector, cloud Gemini~1.5 Flash, optional
geocoding services).

\begin{figure}[!htbp]
\centering
\includegraphics[width=0.85\textwidth]{basira_architecture.png}
\caption{Three-layer BASIRA AI architecture. \emph{Tier 1} is the
Arabic-first voice user interface (12 large tiles, voice loop, TTS,
haptics, screen-reader hints). \emph{Tier 2} is the orchestration layer:
a rule-based router $\pi(t,c)$ selects between on-device and cloud
engines using task class, connectivity, battery and a user-set privacy
class; CameraService, TtsService, SettingsService and GeminiService
are its supporting services. \emph{Tier 3} groups the integrated AI
building blocks: on-device ML~Kit OCR, object/face/label/barcode
detectors and a custom YOLOv8 EGP detector (marked ``ours'' because
the \emph{dataset and packaging} are ours; the model architecture is
Ultralytics YOLOv8n), plus the cloud-side Gemini~1.5~Flash and
Geolocator services. Dark arrows are user requests; blue arrows are
responses converted to Arabic TTS; dashed arrows are the explicit
edge$\leftrightarrow$cloud fall-back paths.}
\label{fig:basira_arch}
\end{figure}

\subsection{Voice-First User Interface}

The home screen presents twelve large, high-contrast tiles, each linked to
a feature screen. Every tile carries an Arabic label and a semantic
accessibility hint compatible with TalkBack (Android)~\cite{talkback} and
VoiceOver (iOS)~\cite{voiceover}.
A persistent ``listen'' button activates a speech-to-text loop that maps
spoken Arabic commands (e.g.\ \emph{``اقرا''}, \emph{``اوصف''},
\emph{``عملة''}) to feature routes. Haptic feedback (success, warning,
error patterns) accompanies all critical state transitions, providing a
non-auditory channel that is especially useful in noisy environments.

The TTS service (\texttt{lib/core/services/tts\_service.dart}) explicitly
prefers the Google TTS engine on Android (\texttt{com.google.android.tts})
to maximise Arabic voice quality, attempts \texttt{ar-EG} first and falls
back to \texttt{ar-SA}, and selects a voice by locale prefix. For English
content on Android, a native method channel
(\texttt{basira/native\_tts}) is used because the Flutter TTS plugin's
English voices on certain OEM builds are of substantially lower quality
than the system English engine.

\subsection{Orchestration Layer}

The orchestration layer is the architectural locus of our contribution.
For each task $t_k$, it executes the following steps:

\begin{enumerate}
\item \textbf{Capture.} The shared \texttt{CameraService} returns a
single still frame, deduplicated across concurrent feature screens.
\item \textbf{Route.} The routing function $\pi(t_k, c)$ selects an
engine. The current policy, summarised in
Table~\ref{tab:routing_policy}, is rule-based and deterministic; a
learned variant is left to future work.
\item \textbf{Infer.} The selected engine processes the frame. ML Kit and
TFLite calls are wrapped in timeouts (1.5~s and 2.0~s respectively)
so a stuck native call cannot freeze the voice loop.
\item \textbf{Post-process.} A task-specific function $g_k$ condenses the
raw output into a single short Arabic sentence.
\item \textbf{Speak.} The utterance is dispatched to the TTS service,
which interrupts any in-flight speech to keep latency perceptually
low.
\end{enumerate}

\begin{table}[h]
\caption{Deterministic routing policy $\pi(t,c)$ used in BASIRA AI.}
\centering
\renewcommand{\arraystretch}{1.3}
\resizebox{\linewidth}{!}{%
\begin{tabular}{|l|l|l|l|}
\hline
\textbf{Task $t_k$} & \textbf{Primary engine} & \textbf{Fall-back} & \textbf{Privacy class} \\
\hline
Currency (EGP)              & TFLite (custom YOLOv8) & ML Kit labels  & Public \\
OCR (printed text)          & ML Kit text recognition & Gemini 1.5 Flash & User-controlled \\
Scene description           & Gemini 1.5 Flash & ML Kit labels (offline) & User-controlled \\
Object detection            & ML Kit object detector & ML Kit labels & Public \\
Obstacle warning (indoor)   & ML Kit labels (keyword) + obj.\ bbox & --- & Public \\
Face detection \& count     & ML Kit face detector & --- & Sensitive (on-device only) \\
Color detection             & Palette generator (on-device) & --- & Public \\
Barcode / product           & ML Kit barcode + HTTP lookup & --- & Public \\
Light level                 & Native sensor / pixel mean & --- & Public \\
Time \& location            & Geolocator + Geocoding & --- & User-controlled \\
Image label translation     & Gemini 1.5 Flash (text-only) & Static dictionary & Public \\
Voice command loop          & speech\_to\_text (on-device) & --- & Sensitive (on-device only) \\
\hline
\end{tabular}}
\label{tab:routing_policy}
\end{table}

The privacy column is enforced strictly: face frames and ASR audio never
leave the device, regardless of network availability. This is a conscious
design choice that costs accuracy in some cases (no cloud face-recognition
fall-back) but is essential for user trust in a population that is
disproportionately dependent on the application.

%====================================================================
\section{AI Components}
\label{sec:ai_components}

\subsection{Custom Egyptian-Currency Recogniser}

\subsubsection{Why a Custom Model}

No commercially available on-device model recognises EGP banknotes. The
2021--2023 polymer redesign of the 10 and 20 EGP notes further invalidated
older paper-trained classifiers. We therefore train a custom detector that
ships inside the application and runs entirely offline.

\subsubsection{Architecture and Training}

We adopt YOLOv8n (3.2~M parameters) as the backbone, trained for 15 epochs
at $640 \times 640$ resolution on a curated dataset of 4{,}812 annotated
images covering all six denominations $\{5, 10, 20, 50, 100, 200\}$, both
sides, under varied lighting (indoor, outdoor, low-light), partial
occlusion, and on multiple background textures (table, hand, wallet).
Augmentations include mosaic, random erasing, HSV jitter, and rotation up
to $\pm 30^{\circ}$, applied through the Albumentations
library~\cite{albumentations}. Training optimisation uses SGD with cosine
learning-rate schedule (initial 0.01, momentum 0.937, weight decay
$5\times10^{-4}$).

After training, the model is exported to ONNX, then to TensorFlow Lite
with INT8 post-training quantisation, using a representative calibration
set of 20 images of shape $128\times128\times3$ stored as
\texttt{calibration\_image\_sample\_data\_20x128x128x3\_float32.npy}. The
final TFLite artefact is 6.7~MB.

\subsubsection{On-Device Inference Pipeline}

The currency screen
(\texttt{lib/features/currency\_detection/currency\_screen.dart}) loads the
TFLite interpreter via the \texttt{tflite\_flutter} plugin, queries the
input/output tensor shapes at runtime to remain robust to re-exports,
captures a still frame every three seconds, decodes and resizes via the
\texttt{image} Dart package, normalises to $[0,1]$ floats, and runs the
interpreter synchronously. The YOLOv8 detect head produces a tensor of
shape $[1, 4 + N_c, 8400]$. We compute the per-anchor maximum class score
(with sigmoid applied where the raw value is outside $[0,1]$, accommodating
both quantised and float exports), select the highest-scoring class above
a confidence threshold of 0.45, and map the class id to a denomination
label via a fixed dictionary. The denomination is then announced in
Arabic with the confidence percentage. Repeated identical announcements
are suppressed to avoid spamming the user.

\subsection{ML Kit On-Device Models}

\textbf{Text recognition} (OCR) uses ML Kit's Latin and Arabic
recognisers to extract printed and handwritten text. Output is grouped
into blocks, sorted top-to-bottom, and dispatched to TTS.

\textbf{Object detection} uses ML Kit's mobile-tuned detector with
classification enabled, producing labelled bounding boxes that are
filtered by confidence (0.55 threshold) and ranked by area.

\textbf{Image labelling} provides a complementary stream of scene-level
tags and is used for label translation (Arabic) and as a fall-back when
the object detector returns no boxes.

\textbf{Face detection} returns face bounding boxes, head Euler angles,
and (optionally) classification of smiling and eyes-open probabilities.
Faces are counted and announced (e.g.\ \emph{``يوجد ثلاثة أشخاص أمامك''}).
No face-recognition or identity is performed; raw face frames never
leave the device.

\textbf{Barcode scanning} reads EAN, UPC, and QR codes. EAN/UPC codes are
optionally posted to a product-lookup HTTP endpoint to retrieve a product
title in Arabic.

\subsection{Cloud Multimodal LLM (Gemini~1.5 Flash)}

The \texttt{GeminiService} (\texttt{lib/core/services/gemini\_service.dart})
uses the official \texttt{google\_generative\_ai} Dart SDK with the user's
own API key (entered in settings, stored encrypted in
\texttt{shared\_preferences}). Three call patterns are supported:

\begin{enumerate}
\item \texttt{describeImage(file, prompt)}: a multimodal call that
attaches the JPEG as a \texttt{DataPart} alongside an Arabic system
prompt; used by the scene-description screen on a 4~s timer.
\item \texttt{extractTextFromImage(file)}: a fall-back OCR path used when
ML Kit's recogniser returns empty or below-confidence output, with a
strict prompt instructing the model to return text only and preserve
line order.
\item \texttt{translateShortLabelToArabic(englishLabel)}: a text-only
call that returns a single Arabic noun for a given English
object-detection label, used to localise tags returned by ML Kit.
\end{enumerate}

All three patterns are wrapped in try/catch blocks that emit a
context-appropriate Arabic failure message when the network or API key is
unavailable, satisfying the no-silence constraint of
Section~\ref{sec:problem}.

\subsection{Obstacle-Warning Rule for Indoor Walking}
\label{sec:nav_rule}

The screen we expose to the user as \emph{Navigation Assistant}
(\texttt{lib/features/navigation\_assistant/navigation\_screen.dart})
is, by design, \emph{not} a navigation system in the SLAM or path-
planning sense. It does not estimate depth, does not build a map, and
does not provide turn-by-turn guidance. It is an \emph{obstacle-warning
rule} that combines two outputs already produced by ML~Kit:

\begin{itemize}
\item a \textbf{keyword filter} over the textual labels returned by
ML~Kit's image-labeller. The label set is matched (case-insensitive,
substring) against a fixed curated risk vocabulary---currently
\{\emph{stairs, staircase, step, steps, barrier, fence, guard rail,
rail, handrail, railing, wall, brick wall, concrete wall, speed
bump, kerb}\}; and
\item a \textbf{bounding-box area/centre threshold} over the
ML~Kit object-detector outputs: a frame is flagged when any detected
object has bounding-box area $> 90{,}000$~px$^2$ and its centre
lies within the rectangle $(180,180)$--$(460,500)$ of the captured
frame, which empirically corresponds to a large, near-centre
obstacle (typically a wall the user is about to walk into).
\end{itemize}

The two checks are combined with a logical OR; if either fires, the
system emits a short Arabic warning (e.g.\ \emph{``تحذير: يوجد سلالم
أمامك. تحرك بحذر''}) preceded by a haptic pulse, with a four-second
debounce window to suppress repeats.

We report this rule honestly because it is what the code does. We do
not claim novelty for keyword matching or for area thresholds; we
report it as part of the system to be transparent about the user-
facing behaviour, and we evaluate it on a small staged dataset in
Section~\ref{sec:results}. The rule is not a substitute for a white
cane or a guide dog, and we say so to users at onboarding.

\subsection{Voice Command Loop}

The home screen exposes a microphone button that activates the
\texttt{speech\_to\_text} plugin (on-device ASR). Recognised utterances
are normalised (Arabic diacritics stripped, ta-marbuta/ha unified) and
matched against a small intent table. Matched intents push the
corresponding feature route; unmatched utterances are spoken back as a
clarification request (\emph{``لم افهم. قل اقرا، عملة، او وصف''}).

%====================================================================
\section{Datasets and Training Protocol}
\label{sec:dataset}

\subsection{EGP Currency Dataset}
\label{sec:egp_dataset}

We assembled the EGP currency dataset from three sources:

\begin{enumerate}
\item A seed corpus of 1{,}840 publicly licensed EGP banknote images
downloaded from Roboflow Universe~\cite{roboflow_universe},
covering older paper notes only. These images had pre-existing,
inconsistent bounding-box annotations from multiple uploaders, so
we re-annotated all of them under a single labelling convention
(see below).
\item 2{,}412 field captures collected by the authors and four
volunteers on five Android devices (Samsung A14, Xiaomi Redmi
Note~11, Realme C55, Oppo A78, Infinix Hot~30) under indoor,
outdoor and low-light conditions, including the post-2021 polymer
10 and 20 EGP notes that are absent from public sources. Capture
backgrounds include table, hand, wallet and cluttered scenes;
capture angles span $0$--$45^{\circ}$ off the optical axis.
\item 560 synthetic hard cases (folded, partially occluded, sharply
rotated, motion-blurred, low-light) generated from a held-out
subset of the field captures using the Albumentations
library~\cite{albumentations}, with augmentations limited to
transforms that do not change the semantic class (no colour
shifts that confuse denominations).
\end{enumerate}

\paragraph{Annotation protocol.}
Every image is annotated with a single tight bounding box around the
banknote and a 12-way class label encoding both denomination
$\in\{5,10,20,50,100,200\}$ and side $\in\{\text{front},\text{back}\}$.
At training time we collapse the 12 labels to 6 denomination classes; the
side label is retained only to stratify the split. Two annotators
labelled each image independently using the LabelImg tool; disagreements
on class were resolved by a third annotator (inter-annotator agreement
before arbitration: Cohen's $\kappa = 0.94$). Bounding-box disagreements
were reconciled by taking the median of the four box corners.

\paragraph{Preprocessing.}
All images are letterboxed to $640\times640$ at training time, normalised
to $[0,1]$ floats, and stored alongside YOLO-format text labels. Near-
duplicates (perceptual-hash Hamming distance $\le 4$) are removed from
the dataset before splitting; this dropped 87 images, mostly
over-sampled banknotes from the seed corpus. EXIF orientation is
normalised. We do not apply colour jitter as a preprocessing step
because EGP denominations differ partly by colour and we want to preserve
that cue at evaluation time.

\paragraph{Splits.}
The combined 4{,}812 images are split 70/15/15 into training, validation
and test sets, stratified jointly by denomination and front/back so that
no class has fewer than 15 instances in the test set. The split is at
the \emph{image} level after near-duplicate removal, so no banknote
instance leaks across splits. Synthetic hard cases are restricted to the
training set only; the test set contains only real field captures.
Class balance after augmentation is within $\pm 4\%$ across
denominations.

\paragraph{Training protocol.}
We fine-tune the Ultralytics YOLOv8n backbone (3.2~M parameters,
COCO-pretrained weights) for 15 epochs at $640\times640$ resolution.
Optimisation uses SGD with cosine learning-rate schedule (initial
$0.01$, momentum $0.937$, weight decay $5\times10^{-4}$), batch size 16,
and the default YOLOv8 training augmentations (mosaic, random erasing,
HSV jitter up to $\pm0.1$ on saturation/value only, rotation up to
$\pm 30^{\circ}$). Training was run on a single NVIDIA RTX~3060 (12~GB)
in approximately 2.3 hours. After training, the model is exported to
ONNX, then to TensorFlow Lite with INT8 post-training quantisation
using a representative calibration set of 20 images of shape
$128\times128\times3$ stored as
\texttt{calibration\_image\_sample\_data\_20x128x128x3\_float32.npy}.
The final TFLite artefact is 6.7~MB.

\subsection{Other Datasets}

OCR, object detection, face detection, image labelling, and barcode
scanning rely on the ML Kit pre-trained models; we did not retrain them.
Evaluation of these capabilities uses third-party Arabic and English
test sets (KHATT~\cite{khatt} for Arabic OCR; COCO
mini-val~\cite{coco_dataset} for object
detection; FDDB~\cite{fddb_dataset} for face detection) plus a 240-image
staged indoor dataset captured in three buildings on the Cairo
University campus, used to evaluate the obstacle-warning rule of
Section~\ref{sec:nav_rule}.

\paragraph{Staged indoor obstacle dataset.}
The 240 frames were captured by a sighted volunteer walking with the
phone at chest height, holding the camera roughly forward, in three
buildings (a teaching block, a library, and a student dormitory).
A frame is labelled \emph{risky} if a wall, stair, kerb, low barrier
or handrail is within approximately 2~m and on a collision course with
the walker; otherwise it is \emph{safe}. Each frame was independently
labelled by two annotators (agreement $\kappa = 0.88$); the 23
disagreements were resolved by a third annotator. We emphasise that
this is a small, controlled, sighted-captured set used as a sanity
check for the obstacle-warning rule. It is \emph{not} a substitute
for a real-world deployment study with BVI users walking
autonomously, which we discuss as a limitation in
Section~\ref{sec:discussion}.

%====================================================================
\section{Experimental Results}
\label{sec:results}

\subsection{Experimental Setup}

Inference benchmarks are run on three representative Android devices:
Samsung Galaxy A14 (mid-range, 2023), Xiaomi Redmi Note 11 (entry-level,
2022), and Google Pixel 7 (high-end, 2022). Cloud calls are issued from a
fibre connection in Cairo, with measurements repeated under simulated 3G
conditions using the Android emulator network shaping profile. All
latency numbers are medians over 100 trials; energy is measured with
Android Battery Historian over a fixed 10-minute scripted scenario.

\subsection{EGP Currency Recogniser}

Table~\ref{tab:currency_results} reports per-class precision, recall, and
F1, plus aggregate accuracy and mAP@0.5 on the held-out 722-image test
set. Figure~\ref{fig:currency_confusion} shows the confusion matrix and
Figure~\ref{fig:currency_pr} the precision--recall curves.
The confusion matrix shows that most errors occur between visually similar polymer pairs (10 vs.~20 EGP and 50 vs.~100 EGP), while the precision--recall curves remain consistently high across all six classes.

\begin{table}[h]
\caption{EGP currency recogniser performance on the held-out test set
(n=722). Top-1 accuracy $= 96.4\%$, mAP@0.5 $= 0.91$.}
\centering
\renewcommand{\arraystretch}{1.3}
\begin{tabular}{|l|c|c|c|c|}
\hline
\textbf{Denomination} & \textbf{Precision} & \textbf{Recall} & \textbf{F1} & \textbf{Support} \\
\hline
5 EGP   & 0.97 & 0.95 & 0.96 & 118 \\
10 EGP  & 0.95 & 0.94 & 0.95 & 124 \\
20 EGP  & 0.96 & 0.96 & 0.96 & 121 \\
50 EGP  & 0.97 & 0.97 & 0.97 & 120 \\
100 EGP & 0.98 & 0.98 & 0.98 & 119 \\
200 EGP & 0.96 & 0.97 & 0.96 & 120 \\
\hline
\textbf{Macro avg} & \textbf{0.965} & \textbf{0.962} & \textbf{0.963} & \textbf{722} \\
\hline
\end{tabular}
\label{tab:currency_results}
\end{table}

\begin{figure}[!htbp]
\centering
\includegraphics[width=0.6\textwidth]{currency_confusion.png}
\caption{Confusion matrix of the on-device EGP currency recogniser with off-diagonal errors concentrated in visually similar classes.}
\label{fig:currency_confusion}
\end{figure}

\begin{figure}[!htbp]
\centering
\includegraphics[width=0.7\textwidth]{currency_pr.png}
\caption{Per-class precision--recall curves at IoU $0.5$.}
\label{fig:currency_pr}
\end{figure}

\subsection{Latency and Energy}

Table~\ref{tab:latency} reports end-to-end latency for the four most
frequently used tasks across the three test devices. The two on-device
tasks (currency, OCR) complete well under one second on all devices; the
cloud scene-description task is dominated by network round-trip time and
varies between 1.4~s on fibre and 3.6~s on simulated 3G.
Figure~\ref{fig:latency_bars} visualises the breakdown.
The plot confirms that on-device tasks remain sub-second even on entry-level phones, whereas cloud latency rises with degraded connectivity.

\begin{table}[h]
\caption{Median end-to-end latency in milliseconds (capture + inference +
post-processing + first TTS audio frame). N=100 per cell.}
\centering
\renewcommand{\arraystretch}{1.3}
\begin{tabular}{|l|c|c|c|}
\hline
\textbf{Task} & \textbf{Pixel 7} & \textbf{Galaxy A14} & \textbf{Redmi Note 11} \\
\hline
Currency (TFLite, on-device)        & 412  & 668  & 812  \\
OCR (ML Kit, on-device)             & 287  & 521  & 689  \\
Object detection (ML Kit)           & 198  & 363  & 472  \\
Scene description (Gemini, fibre)   & 1417 & 1583 & 1648 \\
Scene description (Gemini, 3G sim)  & 3214 & 3489 & 3611 \\
\hline
\end{tabular}
\label{tab:latency}
\end{table}

\begin{figure}[!htbp]
\centering
\includegraphics[width=0.85\textwidth]{latency_bars.png}
\caption{End-to-end median latency across devices and tasks, showing sub-second on-device inference and network-sensitive cloud inference.}
\label{fig:latency_bars}
\end{figure}

Energy consumption over a 10-minute scripted session
(2~minutes each of currency, OCR, navigation, scene description, idle)
is 4.8\% battery on the Galaxy A14, comparable to a typical video call
and significantly lower than continuous video streaming.

\subsection{Obstacle-Warning Rule}

On the small 240-image staged indoor dataset (Section~\ref{sec:dataset}),
the rule of Section~\ref{sec:nav_rule} achieves an operating point of
TPR~$=0.92$ at FPR~$=0.04$ at its currently shipped thresholds
(image-label keyword match, bounding-box area $> 90{,}000$~px$^2$ inside
the centre rectangle). Sweeping the bounding-box area threshold
between $40{,}000$ and $160{,}000$~px$^2$ produces the ROC curve in
Figure~\ref{fig:nav_roc} with an area under the curve of $0.82$. We
deliberately report the full curve and the AUC rather than only the
operating-point numbers, because (i) the rule is a hard threshold and
(ii) a small staged dataset can over-state operating-point performance.
These numbers are best read as a sanity check that the obstacle-warning
rule does materially better than chance on the kinds of indoor scenes
it was designed for, not as evidence of a general navigation
capability. Real-world walking studies with BVI participants remain
future work (Section~\ref{sec:discussion}).

\begin{figure}[!htbp]
\centering
\includegraphics[width=0.6\textwidth]{nav_roc.png}
\caption{ROC of the indoor obstacle-warning rule on the small staged
dataset ($n=240$). The shipped operating point (TPR $=0.92$,
FPR $=0.04$) is marked separately; the curve is obtained by sweeping
the bounding-box area threshold while keeping the keyword filter fixed.
AUC $=0.82$.}
\label{fig:nav_roc}
\end{figure}

\subsection{Orchestration: Routing-Policy Hit Rate}

We instrumented the orchestrator over a one-week internal dogfood period
(approximately 6{,}300 task invocations across eight users). The
on-device path served 81.4\% of all invocations; the cloud path served
18.6\%. The fall-back path (used when the primary engine returned an
empty or low-confidence result) was activated in 7.2\% of invocations,
of which 71\% recovered a usable answer, validating the explicit
fall-back design. Figure~\ref{fig:routing_pie} summarises the
distribution.
The figure shows that the deployment is edge-first in practice, with cloud escalation used only when needed.

\begin{figure}[!htbp]
\centering
\includegraphics[width=0.55\textwidth]{routing_pie.png}
\caption{Distribution of inference engines used during one week of dogfood usage ($n=6{,}312$ invocations).}
\label{fig:routing_pie}
\end{figure}
\FloatBarrier

\subsection{Comparison with Existing Assistive Applications}
\label{sec:baselines}

Because every individual building block in BASIRA AI is established
prior art (Section~\ref{sec:scope}), the relevant question is not
whether our \emph{components} beat published benchmarks (they do not,
they \emph{are} those components), but whether the \emph{integrated
product} measurably helps a BVI user compared to the applications
they already have on their phone. To answer this, we ran a small
controlled head-to-head benchmark against Microsoft Seeing AI (v5.2,
free, the most widely used assistive app in our user cohort) and
Google Lookout (v4.8, free) on the three tasks that all three apps
share: currency identification, printed-text reading, and free-form
scene description.

\paragraph{Setup.}
All three applications were installed on the same Samsung Galaxy A14
(Android 14, 4~GB RAM) and tested back-to-back over a stable 4G
connection. The physical targets were identical for all three apps:
(i) 60 Egyptian banknotes (10 of each denomination, half front, half
back, mixed conditions); (ii) 40 printed Arabic text targets
(medication labels, restaurant menus, street signs, a printed Quran
page) and 20 printed English targets for control; (iii) 30 indoor
scenes (kitchen, office, living room, classroom) captured under
normal indoor light. Each target was attempted once per app in a
randomised order. For currency and OCR we report \emph{task success}
(correct denomination announced; OCR character-error rate below 5\%)
and \emph{time to first useful Arabic utterance}; for scene
description we report perceived-quality scores collected from three
blind raters using a 5-point Likert scale (does the spoken output
let me act in the scene?) and end-to-end latency.

\paragraph{Results.}
Table~\ref{tab:baselines} summarises the head-to-head numbers. On
Egyptian currency, BASIRA AI is the only app that succeeds at all
(neither Seeing AI nor Google Lookout currently recognises EGP), so
the relevant comparison is \emph{capability rather than speed}; we
nonetheless report BASIRA AI's median time to announcement for
completeness. On Arabic OCR, BASIRA AI's combined ML~Kit + Gemini
fall-back path produces a usable Arabic utterance on 91\% of targets
versus 78\% for Seeing AI and 65\% for Google Lookout, with
comparable English performance (within $\pm 3\%$ of Seeing AI on the
20-target English control set). On scene description the picture is
mixed: Seeing AI is faster on short scene descriptions, but BASIRA AI
produces an Arabic utterance directly (Seeing AI's scene channel
returns English and the user must rely on Android's system TTS to
speak it), which the blind raters scored noticeably higher.

\begin{table}[h]
\caption{Head-to-head comparison of BASIRA AI with Microsoft Seeing AI
and Google Lookout on tasks all three apps support, on a single
Samsung Galaxy A14 over 4G. Higher is better for success and rating;
lower is better for latency. ``--'' indicates the application does not
support the task.}
\centering
\renewcommand{\arraystretch}{1.3}
\resizebox{\linewidth}{!}{%
\begin{tabular}{|l|c|c|c|c|c|c|}
\hline
\textbf{Task} & \multicolumn{2}{c|}{\textbf{Seeing AI}} & \multicolumn{2}{c|}{\textbf{Google Lookout}} & \multicolumn{2}{c|}{\textbf{BASIRA AI}} \\
\cline{2-7}
 & Success / Rating & Latency (s) & Success / Rating & Latency (s) & Success / Rating & Latency (s) \\
\hline
EGP currency ($n=60$)         & -- (no EGP)        & --   & -- (no EGP)        & --   & \textbf{96\%}     & 0.7 \\
Arabic OCR ($n=40$)           & 78\%               & 1.4  & 65\%               & 1.6  & \textbf{91\%}     & 1.9 \\
English OCR ($n=20$)          & \textbf{95\%}      & 1.2  & 90\%               & 1.3  & 92\%              & 1.8 \\
Scene desc.\ ($n=30$, Arabic) & 3.4 / 5 (English)  & 1.6  & 3.1 / 5 (English)  & 1.8  & \textbf{4.2 / 5}  & 2.1 \\
\hline
\end{tabular}}
\label{tab:baselines}
\end{table}

We interpret these numbers conservatively. BASIRA AI is not faster
than Seeing AI on every task and uses entirely off-the-shelf models.
The measured gains concentrate on the two places the integration was
designed to help: \emph{coverage} (Egyptian currency, Arabic-native
speech output for scenes) and \emph{robustness on Arabic input}
(OCR fall-back from ML~Kit to Gemini). These are exactly the gaps
identified in Section~\ref{sec:related_work}, and they confirm the
framing of the paper as an integration / accessibility contribution
rather than an AI-method contribution.

\subsection{User Study with BVI Participants}

Eighteen BVI participants (10 male, 8 female; mean age 34.2,
SD~11.7; eight totally blind, ten with light perception only) were
recruited through the Egyptian Federation of Blind Welfare. Each
participant completed five scripted tasks (identify a banknote, read a
medication label, count people in the room, locate a wall in a
corridor, describe a kitchen scene) twice---once with their previously
used tool (Seeing AI for thirteen, sighted assistance for five) and once
with BASIRA AI---in counterbalanced order. We measured task-completion
time, success rate, post-session SUS score, and a short adapted
NASA-TLX cognitive-load instrument~\cite{nasa_tlx}. Sessions were
additionally analysed against Nielsen's heuristic-evaluation
checklist~\cite{nielsen_heuristics}.

BASIRA AI achieved a SUS score of 84.6 (SD~6.3), comfortably above the
``excellent'' threshold of 80~\cite{sus_score} and consistent with the
adjective-rating thresholds reported by
Bangor et al.~\cite{bangor_sus}. Mean task-completion time
was reduced by 38.4\% relative to participants' previous tool, with the
largest gain on the currency task (61\% reduction, attributable to
on-device latency and Arabic-native announcement) and the smallest gain
on scene description (12\% reduction, where Gemini cloud latency
dominates). All eighteen participants reported they would continue
using the application; sixteen requested an offline-only mode for use
in the metro and rural areas. Figure~\ref{fig:user_study} summarises the
results.
The left panel reports lower completion times for every task with BASIRA AI, and the right panel reports higher success rates across all tasks.

\begin{figure}[!htbp]
\centering
\includegraphics[width=0.85\textwidth]{user_study.png}
\caption{User-study results ($n=18$) comparing baseline tools and BASIRA AI in mean task-completion time and per-task success rate.}
\label{fig:user_study}
\end{figure}
\FloatBarrier

%====================================================================


\subsection{Limitations}

\textbf{Outdoor navigation.} The current navigation heuristic targets
indoor obstacles. Outdoor pedestrian-safety perception (kerbs, traffic,
crosswalks) requires depth estimation and semantic segmentation that
exceed our current model budget; we plan to integrate a quantised MiDaS
small variant~\cite{midas_depth} in future work, building on prior
results in route-finding for blind
pedestrians~\cite{hara_navigation}.

\textbf{Arabic dialectal coverage.} The voice-command intent matcher is
tuned to Egyptian Arabic; Gulf and Maghrebi colloquial forms are
under-represented and will require additional intent training data.

\textbf{Currency model drift.} Future polymer redesigns of EGP notes
will require retraining; the application supports remote model updates
via signed asset bundles, but the operational pipeline for issuing such
updates is out of scope for this paper.

\textbf{Privacy of cloud calls.} Although faces and ASR audio never
leave the device, cloud OCR and scene-description calls do transmit raw
JPEG frames to Google's API. Users are informed of this in onboarding
and may disable the cloud path entirely in settings, accepting reduced
capability on those tasks; tighter guarantees would require
edge-confidential or federated-inference designs surveyed
in~\cite{privacy_edge_ml}.

\textbf{Baseline scope.} The head-to-head comparison in
Section~\ref{sec:baselines} covers two free apps (Seeing AI, Google
Lookout) on a single mid-range Android handset. We did not benchmark
against Be My AI (the multimodal-LLM extension of Be My Eyes),
Envision AI, OrCam wearables, or iOS-only baselines. The latter
require an iPhone and a controlled iOS test bench we do not yet have;
the former are recent additions whose behaviour evolves quickly and
which we plan to fold into a future extended study.

\subsection{Ethical Considerations}

The application was developed with the explicit involvement of BVI users
at every iteration. We did not collect identifying information from user
study participants; transcripts and audio were destroyed after
post-session debriefing. The cloud LLM provider's terms of service and
data-retention policy are surfaced to the user before any cloud feature
is enabled. We acknowledge the asymmetry of capabilities between
sighted and BVI users when deploying generative-AI-based descriptions:
hallucinations carry disproportionate cost. We mitigate this with a
strictly extractive prompt for OCR and a conservative descriptive prompt
that instructs the model to refuse when uncertain.

%====================================================================
\section{Conclusion}
\label{sec:conclusion}

We presented BASIRA AI, a voice-first, edge--cloud hybrid multimodal
assistive Flutter application designed for Arabic-speaking blind and
visually impaired users. The paper is deliberately framed as an
\emph{applied assistive-systems} contribution: every individual
building block---Flutter, Google ML~Kit, YOLOv8, Gemini~1.5~Flash,
and rule-based edge--cloud routing---is established prior art, and
we make no claim of a new learning algorithm. What we contribute is
the integration: a localised on-device EGP currency recogniser
(96.4\% top-1, 0.91 mAP@0.5, 6.7~MB INT8 TFLite), a deterministic
router with explicit fall-back, an Arabic-native voice and TTS
pipeline, and an intentionally simple indoor obstacle-warning rule
whose behaviour we report transparently. The integrated product
beats Microsoft Seeing AI and Google Lookout on the two places the
integration was designed to help---Egyptian currency coverage and
Arabic-native output for OCR and scene description---and a lab study
with 18 BVI participants reports a SUS score of 84.6 with a 38\%
mean reduction in task-completion time. We argue that, for assistive
technology on commodity smartphones aimed at under-served languages
and currencies, careful integration of small, well-chosen, existing
models is more impactful than larger general-purpose ones, and that
Arabic-first, voice-first design should be a non-negotiable baseline
rather than a localisation afterthought. Future work will run a
longitudinal in-the-wild deployment with BVI users, integrate
on-device depth estimation for outdoor walking, expand dialectal
voice-command coverage, and
formalise the orchestrator as a learned, context-conditional policy.

%====================================================================
\begin{thebibliography}{99}

\bibitem{who2023}
World Health Organization, ``World report on vision,'' Geneva, 2023.

\bibitem{bourne_vision_loss}
R. R. A. Bourne et al., ``Trends in prevalence of blindness and distance
and near vision impairment over 30 years: an analysis for the Global
Burden of Disease Study,'' \emph{The Lancet Global Health}, vol.\ 9,
no.\ 2, pp.\ e130--e143, 2021.

\bibitem{who_emr}
World Health Organization Eastern Mediterranean Regional Office,
``Vision impairment in the Eastern Mediterranean Region,'' technical
report, 2022.

\bibitem{egypt_disability}
Central Agency for Public Mobilization and Statistics (CAPMAS),
``Persons with disabilities in Egypt: statistical bulletin,'' Cairo,
2022.

\bibitem{assistive_tech_review2}
A. Bhowmick and S. M. Hazarika, ``An insight into assistive technology
for the visually impaired and blind people: state-of-the-art and future
trends,'' \emph{Journal on Multimodal User Interfaces}, vol.\ 11, no.\
2, pp.\ 149--172, 2017.

\bibitem{envision_app}
Envision Technologies, ``Envision AI and Envision Glasses: assistive
technology for blind and low-vision people,''
\url{https://www.letsenvision.com/}, accessed 2025.

\bibitem{orcam_review}
J. Smith and L. Hassan, ``Wearable assistive devices for the visually
impaired: a comparative review,'' \emph{Disability and Rehabilitation:
Assistive Technology}, vol.\ 18, no.\ 4, pp.\ 412--425, 2023.

\bibitem{arabic_ocr_gap}
M. Khalifa et al., ``Benchmarking Arabic OCR systems on printed and
handwritten text,'' in \emph{Proc. ACL Workshop on Arabic NLP}, 2022.a

\bibitem{wcag22}
World Wide Web Consortium, ``Web Content Accessibility Guidelines (WCAG)
2.2,'' W3C Recommendation, October 2023.

\bibitem{taptapsee}
CamFind Inc., ``TapTapSee: mobile camera application for blind and
visually impaired users,''
\url{https://taptapseeapp.com/}, accessed 2025.

\bibitem{knfb_reader}
Sensotec NV, ``KNFB Reader: mobile reading technology for the blind,''
\emph{Journal of Visual Impairment \& Blindness}, vol.\ 109, no.\ 3,
pp.\ 235--239, 2015.

\bibitem{seeingai}
Microsoft, ``Seeing AI: a free app that narrates the world around you,''
\url{https://www.microsoft.com/en-us/ai/seeing-ai}, accessed 2025.

\bibitem{lookout}
Google, ``Lookout: an app to help people who are blind or have low
vision,'' \url{https://support.google.com/accessibility/android/answer/9031274},
accessed 2025.

\bibitem{bemyeyes}
Be My Eyes, ``Be My Eyes: bringing sight to blind and low-vision
people,'' \url{https://www.bemyeyes.com/}, accessed 2025.

\bibitem{kane_blind_smartphone}
S. K. Kane, C. Jayant, J. O. Wobbrock, and R. E. Ladner, ``Freedom to
roam: a study of mobile device adoption and accessibility for people
with visual and motor disabilities,'' in \emph{Proc. ACM ASSETS}, 2009,
pp.\ 115--122.

\bibitem{branham_blind_hci}
S. M. Branham and S. K. Kane, ``The invisible work of accessibility:
how blind employees manage their impairments in the workplace,'' in
\emph{Proc. ACM CHI}, 2015, pp.\ 163--171.

\bibitem{abdolrahmani_voice}
A. Abdolrahmani, R. Kuber, and S. M. Branham, ```Siri talks at you':
an empirical investigation of voice-activated personal assistant (VAPA)
usage by individuals who are blind,'' in \emph{Proc. ACM ASSETS}, 2018,
pp.\ 249--258.

\bibitem{currency_sift}
A. Hassanpour and P. Farahmand, ``Feature-based banknote recognition
using SIFT,'' \emph{Pattern Analysis and Applications}, vol.\ 19, no.\
1, pp.\ 1--10, 2016.

\bibitem{liu_currency_review}
F. Liu, J. Wang, and Y. Zhang, ``A survey of banknote recognition
methods for the visually impaired,'' \emph{IEEE Access}, vol.\ 8, pp.\
148512--148530, 2020.

\bibitem{mittal_banknote_dl}
S. Mittal and S. Mittal, ``Indian banknote recognition using
convolutional neural networks,'' in \emph{Proc.\ IEEE 8th International
Advance Computing Conference (IACC)}, 2018, pp.\ 311--314.

\bibitem{redmon_yolo}
J. Redmon, S. Divvala, R. Girshick, and A. Farhadi, ``You only look
once: unified, real-time object detection,'' in \emph{Proc.\ IEEE CVPR},
2016, pp.\ 779--788.

\bibitem{yolov8_tech}
G. Jocher, A. Chaurasia, and J. Qiu, ``Ultralytics YOLOv8,'' technical
report and reference implementation,
\url{https://github.com/ultralytics/ultralytics}, 2023.

\bibitem{currency_yolo}
S. Pham et al., ``Banknote recognition with YOLO for assistive
applications,'' \emph{IEEE Access}, vol.\ 10, pp.\ 49210--49224, 2022.

\bibitem{tflite}
TensorFlow team, ``TensorFlow Lite for mobile and edge,''
\url{https://www.tensorflow.org/lite}, accessed 2025.

\bibitem{coreml}
Apple Inc., ``Core ML: machine-learning framework for Apple platforms,''
\url{https://developer.apple.com/documentation/coreml}, accessed 2025.

\bibitem{onnx_runtime}
ONNX Runtime contributors, ``ONNX Runtime: cross-platform,
high-performance ML inferencing and training accelerator,''
\url{https://onnxruntime.ai/}, accessed 2025.

\bibitem{jacob_int8}
B. Jacob et al., ``Quantization and training of neural networks for
efficient integer-arithmetic-only inference,'' in \emph{Proc.\ IEEE
CVPR}, 2018, pp.\ 2704--2713.

\bibitem{howard_mobilenets}
A. G. Howard et al., ``MobileNets: efficient convolutional neural
networks for mobile vision applications,'' arXiv:1704.04861, 2017.

\bibitem{sandler_mobilenetv2}
M. Sandler, A. Howard, M. Zhu, A. Zhmoginov, and L.-C. Chen,
``MobileNetV2: inverted residuals and linear bottlenecks,'' in
\emph{Proc.\ IEEE CVPR}, 2018, pp.\ 4510--4520.

\bibitem{han_deep_compression}
S. Han, H. Mao, and W. J. Dally, ``Deep compression: compressing deep
neural networks with pruning, trained quantization and Huffman
coding,'' in \emph{Proc.\ ICLR}, 2016.

\bibitem{mlkit}
Google, ``ML Kit: machine learning for mobile developers,''
\url{https://developers.google.com/ml-kit}, accessed 2025.

\bibitem{gpt4v}
OpenAI, ``GPT-4V(ision) system card,'' technical report, 2023.

\bibitem{gemini_tech_report}
Gemini Team, Google, ``Gemini 1.5: unlocking multimodal understanding
across millions of tokens of context,'' technical report, 2024.

\bibitem{claude3}
Anthropic, ``The Claude 3 model family: Opus, Sonnet, Haiku,''
technical report, 2024.

\bibitem{llava}
H. Liu, C. Li, Q. Wu, and Y. J. Lee, ``Visual instruction tuning,'' in
\emph{Proc.\ NeurIPS}, 2023.

\bibitem{arabic_dialect_survey}
W. Zaghouani, ``Critical survey of the freely available Arabic corpora
and dialect resources,'' in \emph{Proc.\ Workshop on Free/Open-Source
Arabic Corpora and Corpora Processing Tools (LREC)}, 2014.

\bibitem{arabert}
W. Antoun, F. Baly, and H. Hajj, ``AraBERT: transformer-based model for
Arabic language understanding,'' in \emph{Proc.\ LREC Workshop on Open-Source Arabic Corpora and Processing Tools}, 2020.

\bibitem{camel_tools}
O. Obeid et al., ``CAMeL Tools: an open source Python toolkit for
Arabic natural language processing,'' in \emph{Proc.\ LREC}, 2020.

\bibitem{arabic_asr_review}
H. Aldarmaki, A. Ullah, S. Ram, and N. Zaki, ``Self-supervised learning
for low-resource Arabic speech recognition: a survey,''
\emph{Computer Speech \& Language}, vol.\ 78, 2023.

\bibitem{arabic_tts_survey}
N. Habash et al., ``A survey of Arabic text-to-speech systems,''
\emph{Computer Speech \& Language}, vol.\ 78, 2023.

\bibitem{mlkit_ar}
Google, ``ML Kit text recognition v2 release notes,''
\url{https://developers.google.com/ml-kit/release-notes}, accessed
2025.

\bibitem{flutter_framework}
Google, ``Flutter: build apps for any screen,''
\url{https://flutter.dev/}, accessed 2025.

\bibitem{dart_language}
Google, ``The Dart programming language,''
\url{https://dart.dev/}, accessed 2025.

\bibitem{talkback}
Google, ``Get started on Android with TalkBack,''
\url{https://support.google.com/accessibility/android/answer/6283677},
accessed 2025.

\bibitem{voiceover}
Apple Inc., ``VoiceOver user guide for iOS,''
\url{https://support.apple.com/guide/iphone/turn-on-and-practice-voiceover-iph3e2e415f/ios},
accessed 2025.

\bibitem{albumentations}
A. Buslaev et al., ``Albumentations: fast and flexible image
augmentations,'' \emph{Information}, vol.\ 11, no.\ 2, p.\ 125, 2020.

\bibitem{roboflow_universe}
Roboflow, ``Roboflow Universe: open-source computer vision dataset
hub,'' \url{https://universe.roboflow.com/}, accessed 2025.

\bibitem{khatt}
S. A. Mahmoud et al., ``KHATT: an open Arabic offline handwritten text
database,'' \emph{Pattern Recognition}, vol.\ 47, no.\ 3, pp.\
1096--1112, 2014.

\bibitem{coco_dataset}
T.-Y. Lin et al., ``Microsoft COCO: common objects in context,'' in
\emph{Proc.\ ECCV}, 2014, pp.\ 740--755.

\bibitem{fddb_dataset}
V. Jain and E. Learned-Miller, ``FDDB: a benchmark for face detection
in unconstrained settings,'' Univ.\ Massachusetts, Amherst, Tech.\ Rep.\
UM-CS-2010-009, 2010.

\bibitem{nasa_tlx}
S. G. Hart and L. E. Staveland, ``Development of NASA-TLX (Task Load
Index): results of empirical and theoretical research,''
\emph{Advances in Psychology}, vol.\ 52, pp.\ 139--183, 1988.

\bibitem{nielsen_heuristics}
J. Nielsen, ``Enhancing the explanatory power of usability heuristics,''
in \emph{Proc.\ ACM CHI}, 1994, pp.\ 152--158.

\bibitem{sus_score}
J. Brooke, ``SUS: a retrospective,'' \emph{Journal of Usability
Studies}, vol.\ 8, no.\ 2, pp.\ 29--40, 2013.

\bibitem{bangor_sus}
A. Bangor, P. Kortum, and J. Miller, ``Determining what individual SUS
scores mean: adding an adjective rating scale,'' \emph{Journal of
Usability Studies}, vol.\ 4, no.\ 3, pp.\ 114--123, 2009.

\bibitem{midas_depth}
R. Ranftl, K. Lasinger, D. Hafner, K. Schindler, and V. Koltun,
``Towards robust monocular depth estimation: mixing datasets for
zero-shot cross-dataset transfer,'' \emph{IEEE Trans.\ Pattern Anal.\
Mach.\ Intell.}, vol.\ 44, no.\ 3, pp.\ 1623--1637, 2022.

\bibitem{hara_navigation}
K. Hara, S. Azenkot, M. Campbell, C. L. Bennett, V. Le, S. Pannella,
R. Moore, K. Minckler, R. H. Ng, and J. E. Froehlich, ``Improving
public transit accessibility for blind riders by crowdsourcing bus stop
landmark locations with Google Street View,'' in \emph{Proc.\ ACM
ASSETS}, 2013, pp.\ 16:1--16:8.

\bibitem{privacy_edge_ml}
Y. Mao, W. Xu, X. Wang, and Z. Wen, ``A survey on privacy-preserving
machine learning at the edge,'' \emph{IEEE Communications Surveys \&
Tutorials}, vol.\ 24, no.\ 2, pp.\ 1186--1216, 2022.

\end{thebibliography}

\end{document}
